%!TEX program = xelatex
\documentclass[lang=en,11pt,a4paper,citestyle =authoryear]{elegantpaper}

% 标题
\title{Homework08 - MATH 765}
\author{Boren(Wells) Guan}

% 本文档命令
\usepackage{array,url,stix}
\usepackage{subfigure}
\usepackage{tikz}
\usepackage{tikz-cd}
\usepackage{mathtools}
\newcommand{\ccr}[1]{\makecell{{\color{#1}\rule{1cm}{1cm}}}}
\newcommand{\code}[1]{\lstinline{#1}}
\newcommand{\prvd}{$\hfill \qedsymbol$}
\newcommand{\Z}{\mathbb{Z}}
\newcommand{\R}{\mathbb{R}}
\newcommand{\N}{\mathbb{N}}
\newcommand{\C}{\mathbb{C}}
\newcommand{\Q}{\mathbb{Q}}
\newcommand{\M}{\mathcal{M}}
\newcommand{\B}{\mathcal{B}}
\newcommand{\X}{\mathcal{X}}
\newcommand{\Hil}{\mathcal{H}}
\newcommand{\range}{\mathcal{R}}
\newcommand{\nul}{\mathcal{N}}
\newcommand{\ip}[2]{\left\langle #1,#2\right\rangle}
\newcommand{\Rc}{\operatorname{Rc}}
\newcommand{\Rm}{\operatorname{Rm}}
\newcommand{\Ric}{\operatorname{Ric}}
\newcommand{\Scal}{\operatorname{S}}
\newcommand{\del}{\partial}

% 文档区
\begin{document}

% 标题
\maketitle

\subsection*{Before Reading:}\par
Since the assignment only asks us to submit three problems, I decided to write complete solutions for all six problems first, and then add a short evaluation at the end. In this way I can compare them before deciding which three are the best to turn in.\par
I will follow the notation used in class as much as possible. In particular, I write the Riemann tensor as $R_{ijk\ell}$ and use the sign convention compatible with Problem 1, namely the hyperbolic space should satisfy
\[
R_{ijk\ell}=-(g_{ik}g_{j\ell}-g_{i\ell}g_{jk}).
\]
\subsection*{Ex.2}
Show that the Ricci and scalar curvatures of the sphere $g=g_{S^n}$ are given respectively by
\[
\Rc=(n-1)g,
\qquad
S=n(n-1).
\]
\textbf{Sol.} \par
The standard sphere $S^n$ has constant sectional curvature $K\equiv 1$. Therefore its Riemann tensor is
\[
R_{ijk\ell}=g_{ik}g_{j\ell}-g_{i\ell}g_{jk}.
\]
If we would like to compute the Ricci tensor, we contract the first and third indices:
\[
R_{j\ell}=g^{ik}R_{ijk\ell},
\]
and hence
\[
R_{j\ell}=g^{ik}(g_{ik}g_{j\ell}-g_{i\ell}g_{jk}).
\]
Then
\[
g^{ik}g_{ik}=n,
\qquad
g^{ik}g_{i\ell}g_{jk}=\delta^k_{\ell}g_{jk}=g_{j\ell},
\]
and
\[
R_{j\ell}=ng_{j\ell}-g_{j\ell}=(n-1)g_{j\ell},
\]
which means
\[
\Rc=(n-1)g.
\]
Now we take one more trace to get the scalar curvature:
\[
S=g^{j\ell}R_{j\ell}=g^{j\ell}(n-1)g_{j\ell}=(n-1)n.
\]
Therefore $S=n(n-1)$ and which is exactly the general rule for a space form of constant sectional curvature $K$: one has
\[
\Rc=(n-1)Kg,
\qquad
S=n(n-1)K.
\]
For the sphere, $K=1$, so we obtain the stated formulas.\par

\subsection*{Ex.3}
Show that the sectional curvature at any point of a Riemannian $2$-manifold is equal to $S/2$.\par
\textbf{Sol.} \par
Let $(M^2,g)$ be a Riemannian $2$-manifold and let $p\in M$ with an orthonormal basis $\{e_1,e_2\}$ for $T_pM$. Since the tangent space is $2$-dimensional, the only $2$-plane is the whole tangent space itself, so the sectional curvature at $p$ is
\[
K(p)=R_{1212}.
\]
where we use the orthonormal basis, so the denominator in the definition of sectional curvature is $1$.

Now let us compute the scalar curvature. By definition,
\[
S=g^{ik}g^{j\ell}R_{ijk\ell}.
\]
For an orthonormal basis, $g^{ij}=\delta^{ij}$, so
\[
S=\sum_{i,j=1}^2 R_{ijij}.
\]
By the symmetries of the Riemann tensor,
\[
R_{1111}=R_{2222}=0,
\qquad
R_{1221}=-R_{1212},
\qquad
R_{2112}=-R_{2121},
\qquad
R_{2121}=R_{1212},
\]
and hence the only nonzero contributions is
\[
S=R_{1212}+R_{2121}=2R_{1212}.
\]
Since $R_{1212}=K$, we get
\(
S=2K
\)
and then
\(
K=\frac{S}{2}
\). So on a $2$-dimensional Riemannian manifold, the scalar curvature contains exactly the same information as the sectional curvature.\par

\subsection*{Ex.6}
A (Riemannian) Einstein metric is one which satisfies
\[
\Rc(g)=\lambda g
\]
for a smooth function $\lambda=\lambda(x)$ on $M$.
\begin{itemize}
    \item[(a)] Show that $\lambda=\dfrac{S}{n}$.
    \item[(b)] Show that if $n\ge 3$ then $S$ is locally constant.
\end{itemize}
\textbf{Sol.} \par
(a) By assumption, we have
\(
R_{ij}=\lambda g_{ij}
\) and then let us take the trace with $g^{ij}$:
\[
S=g^{ij}R_{ij}=g^{ij}(\lambda g_{ij})=\lambda g^{ij}g_{ij}=\lambda n,
\]
which means
\(
\lambda=\frac{S}{n}
\)
.

(b) Now use the contracted Bianchi identity
\[
\nabla_i R_{ij}=\frac12\nabla_j S.
\]
Since $R_{ij}=\lambda g_{ij}$, metric compatibility $\nabla g=0$ gives
\[
\nabla_i R_{ij}=\nabla_i(\lambda g_{ij})=(\nabla_i\lambda)g_{ij}+\lambda\nabla_i g_{ij}=\nabla_j\lambda.
\]
Then we know
\(
\nabla_j\lambda=\frac12\nabla_j S.
\)
and by (a), $\lambda=S/n$, so
\(
\nabla_j\lambda=\frac1n\nabla_j S
\).

Comparing the two expressions for $\nabla_j\lambda$, we get
\(
\frac1n\nabla_j S=\frac12\nabla_j S
\) and hence
\[
\left(\frac1n-\frac12\right)\nabla_j S=0.
\]
If $n\ge 3$, we have\(
\nabla_j S=0
\)
for every $j$. So $S$ is constant on each connected component, equivalently locally constant. Therefore
\(
\text{if }n\ge 3,\text{ then }S\text{ is locally constant}.
\)
In particular, when $n\ge 3$, an Einstein metric satisfies
\[
\Rc=\frac{S}{n}g
\]
with constant coefficient on each connected component.\par


\addappheadtotoc

\end{document}
